% Chapter 2, section 3: Push messages
% Write in plain LaTeX. Code of a guide step: \guide{<step>}; commands: \begin{shell}...\end{shell};
% a screenshot: \screenshot{file}{caption}; a gap in the draft: \placeholder{...}.
% See book/latex/queen.sty for the full list.

\section{Push messages}

In Kafka you choose a number of partitions and hash keys onto them, so a slow customer holds up
every customer that shares its partition. Here a partition is one entity: a customer, an order, a
conversation. The push that names it creates it, order holds inside it, and one queue has held ten
million.

The \texttt{transactionId} makes a retry safe: the same id within the dedup window (an hour by default) is
not written twice.

\guide{push}

\why{order matters per entity (an order's events must be handled in sequence), not across the
whole queue. One partition per entity keeps that order without one slow entity holding up the
others. A push with no partition goes to a partition called \texttt{Default}.}

\subsection{Every message gets its own answer}

One push can carry many messages, and the broker answers each one: \texttt{queued} (stored, with
its offset), \texttt{duplicate} (that \texttt{transactionId} is already there, nothing was written)
or \texttt{error} (refused, nothing written). The JavaScript client throws when any message is
refused, unless you handle it yourself, and never throws for a duplicate.

\queenfig{push-verdicts}{One push, one answer per message.}

\guide{push-callbacks}

\why{a duplicate is usually good news: it means an earlier attempt already got through. A refusal
is not. Telling them apart per message lets one bad item in a batch of a hundred be retried alone,
and lets a producer that retried after a timeout know that its first attempt worked.}

\subsection{Buffering, for many small messages}

Each push is one HTTP request. A producer that pushes thousands of small messages a second can
let the client collect them and send them in batches instead.

\guide{push-buffer}

A buffer is sent when it holds \texttt{messageCount} messages or when \texttt{timeMillis} has
passed, whichever comes first. If the producer gets far ahead of the network (four times
\texttt{messageCount} waiting), a buffered push waits instead of letting memory grow.

\queenfig{push-buffer}{Client buffering: many small pushes, a few requests.}

\why{one request per message spends most of its time on the round trip, not on the message.
Batching gets far more messages through the same connection. The price: a message is stored only
when its batch is sent, so a process that crashes loses what was still in its buffer. Call
\texttt{flushAllBuffers()} before you exit, and do not buffer messages you cannot afford to lose.}

\subsection{When several messages must land together}

A push of several messages is not all-or-nothing: each message gets its own answer, so some can be
stored while others are refused. When messages must be stored together or not at all (an order and
its payment request, say), push them in a transaction (see Advanced transactional usage).

\subsection{Binary push, compression and framed pop}

By default the client and the broker speak JSON. Three options of the client change how messages
travel, and nothing else: the calls, their results, deduplication and the guarantees stay the
same.

\begin{options}{Option}{Default}{What it does}
binaryPush & false & A push, direct or buffered, goes as one binary body instead of a JSON document. A transaction that pushes sends its pushes the same way. \\
pushCompression & false & Needs \texttt{binaryPush}. The client packs each partition's batch the way the broker stores it, and compresses it with zstd once. The broker keeps those bytes as they are. \\
framedPop & false & \texttt{pop()} and \texttt{consume()} get the records as they are stored, and the client unpacks them into the same messages. \\
\end{options}

Turn them on when you create the client:

\guide{wire-forms}

zstd comes with Node.js, so there is nothing to install. Some requests stay JSON by themselves: a
push with a \texttt{traceId}, a pop of a conflating group, and a pop of an encrypted queue. Acks,
KV, timers and locks do not change.

\why{at a high message rate, the broker spends CPU turning JSON into stored records
and back, and the network carries every message uncompressed. These options move that work into
the clients, where many processes share it, and send the batches compressed. When neither the
broker's CPU nor the network is the limit, leave them off: JSON is easier to read when you debug.}

One trap: \texttt{binaryPush} and \texttt{pushCompression} need broker 2.3.0 or later on every node.
An older broker answers 400, and the client does not fall back to JSON, so turn them on after the
last node is upgraded. \texttt{framedPop} is safe with any broker: an older one ignores it and
answers JSON.

On the site, with the code in all six SDK languages: \href{https://queenmq.com/writing-applications/push/}{queenmq.com/writing-applications/push}. More there: \href{https://queenmq.com/writing-applications/clients/}{the binary forms in every SDK}.
